\section{Related Work and Method Selection}
\label{sec:related}

\subsection{Approaches to industrial anomaly detection}
\emph{Reconstruction} methods train an autoencoder, VAE or GAN on normal images and use the reconstruction error as the anomaly map~\cite{bergmann2019ssim,akcay2018ganomaly,kingma2014vae}.
These networks tend to generalise too well: they reconstruct defects too and blur fine textures, so they lag behind feature-based methods on MVTec AD~\cite{bergmann2021mvtec}.
\emph{Feature-embedding} methods instead model the distribution of pretrained deep features.
A single Gaussian over global ImageNet features is already a strong image-level detector~\cite{rippel2020gaussian}.
SPADE~\cite{cohen2020spade}, PaDiM~\cite{defard2021padim} and PatchCore~\cite{roth2022patchcore} extend the idea to patches, which yields localisation.
\emph{Teacher--student} methods train a student to regress a frozen teacher's features on normal data and read anomalies from the discrepancy~\cite{bergmann2020students,wang2021stfpm,deng2022rd4ad,batzner2024efficientad}.
\emph{Normalising flows} learn an exact likelihood over pretrained features~\cite{rudolph2021differnet,gudovskiy2022cflow,yu2021fastflow}.
\emph{Self-supervised} methods synthesise pseudo-anomalies and train a discriminator~\cite{li2021cutpaste,zavrtanik2021draem,liu2023simplenet}.
\emph{Foundation-model} approaches use CLIP~\cite{radford2021clip} or DINOv2~\cite{oquab2024dinov2} features for zero- or few-shot detection~\cite{jeong2023winclip,damm2025anomalydino}.
\emph{Unified} models train one network for all categories~\cite{you2022uniad}.
The last five families require training, and the zero-shot and unified models still trail the best per-category methods in accuracy.

\subsection{Selected methods}
We select two feature-embedding methods that share a frozen ImageNet WideResNet-50-2~\cite{zagoruyko2016wrn,deng2009imagenet} backbone but model normality in different ways.
\textbf{PaDiM} (baseline) is \emph{parametric} and \emph{position-specific}: it fits one Gaussian per spatial position and therefore assumes that a given grid cell shows the same part of the product in every image.
\textbf{PatchCore} (main method) is \emph{non-parametric} and \emph{position-agnostic}: a test patch only needs to resemble \emph{some} normal patch from any image at any location, and the memory bank covers multi-modal normal appearance.
Four considerations motivate this choice.
\emph{Accuracy}: PatchCore is among the strongest full-shot methods on MVTec AD (99.1\% image AUROC, 93.4\% AU-PRO~\cite{roth2022patchcore}), and PaDiM (95.3\%, 92.1\%~\cite{defard2021padim}) is a clearly weaker baseline, so the gap between them is large enough to analyse.
\emph{Complexity}: neither method uses back-propagation, and both have few interpretable hyper-parameters, which keeps failure analysis tractable.
\emph{Compute}: fitting is a single forward pass followed by covariance estimation or coreset selection, and takes seconds on one GPU.
\emph{Applicability}: both need only normal images, and PatchCore's memory bank can be extended without retraining and explains each decision by its nearest normal patch.
Because the backbone is shared, any performance difference reflects the modelling principle.
We expected the gap to be largest on screw and capsule, where pose and shift variation would violate PaDiM's alignment assumption, and expected the stationary textures to suit both methods; Section~\ref{sec:results} tests this hypothesis.
